\section{Results}

\subsection{Sample and analytic strategy}

The analysis sample was the harmonized baseline (visit V0) of the three FACE
cohorts: $N=9{,}013$ patients (bipolar disorder [BP] $6{,}252$; schizophrenia [SZ]
$2{,}209$; major depression [DR] $552$) recruited across $21$ expert-centre sites
(Table~\ref{tab:characteristics}). The cohorts differ in age, sex distribution and
chronicity, as expected for these disorders. We deliberately fit the measurement
model on the \emph{full} sample without completeness selection: $143$ indicators were
modelled, with a mean cell missingness of $39.8\%$ preserved as missing and never
imputed. Diagnosis (cohort and DSM-5 subtype) was held out of the model and used only
as a validation or metadata label. The reported map is one global,
cohort-reweighted, full-sample fit estimated with a Gaussian-copula
(rank-based) marginal model; because the same structure is recovered under a plain
Gaussian likelihood, the central findings do not depend on that modelling choice---a
robustness result, not an assumption (Annex~\ref{ann:measurement}). Follow-up visits
at $12$ and $24$ months (retained: $47.4\%$ and $32.8\%$ of V0; \edfig{fig:consort})
were used only to test temporal coherence and prognosis, never to define the map. The
high-level model is summarized below and in Methods; its complete derivation,
including the observed-data likelihood, the rotation argument, the
covariance identity and the equivalence with full-information maximum likelihood, is
given in Annex~\ref{ann:measurement}. Figure~\ref{fig:glance} presents the atlas at a
glance; the study design and the three load-bearing invariants are detailed in
\edfig{fig:design}.

\begin{figure}[t]
  \centering
  \IfFileExists{figures/fig1_glance.png}{%
    \includegraphics[width=\linewidth]{fig1_glance.png}%
  }{%
    \fbox{\parbox[c][6.6cm][c]{0.96\linewidth}{\centering\itshape
      Figure~1 (graphical abstract) renders here once the overview infographic is
      saved to \texttt{article/figures/fig1\_glance.png}.}}%
  }
  \caption{\textbf{A latent atlas of severe mental illness, at a glance.} The global,
  missingness-aware Bayesian bifactor / exploratory structural-equation model, fit to
  the FACE baseline ($N=9{,}013$) with diagnosis held out, yields a nine-dimension
  transdiagnostic latent space---a general functional-burden factor (\Gfac) and eight
  specific dimensions. The patient cloud is a graded continuum (best silhouette $0.140$,
  indistinguishable from a single-Gaussian null $0.141\pm0.002$), read through four
  archetypal extremes whose convex blends describe every patient: a \emph{biological}
  corner (A0; high metabolic, inflammatory and substance load), a \emph{low-burden} pole
  (A1), a \emph{severe, non-biological} corner (A2), and a \emph{symptom-driven} corner
  (A3). Metabolic and inflammatory load are largely dissociated from general burden
  (correlation $0.12$ and $0.07$, versus $0.39$ for cognition and $0.42$ for sleep, and
  robust to medication, adiposity and site); biology is durable (metabolic test--retest
  intraclass correlation $0.91$); the archetype phenotype improves held-out two-year
  functional prediction ($\Delta$ELPD $+59.4\pm11.1$); and two-year functional remission
  ranges $27\%$--$60\%$ across corners, the biological corner worst. A0 and A2 carry
  comparable severity but opposite biology.}
  \label{fig:glance}
\end{figure}

\subsection{A nine-dimension transdiagnostic map (the map exists)}

From ten candidate constructs, one global model returned a stable nine-dimension
structure: a general functional-burden factor (\Gfac) plus eight weakly correlated
specific axes---cognition, metabolic load, inflammation, sleep, developmental risk,
suicidality, mania and substance use (Figure~\ref{fig:atlas};
Table~\ref{tab:dimensions}). Inter-factor correlations among the specific axes were
weak throughout (mean absolute off-diagonal $\approx 0.10$), confirming genuinely
distinct axes rather than one collapsed dimension. \Gfac{} was anchored only by
functioning and global-severity indicators (FAST loading $0.90$, EGF $0.73$) and
carried no symptom content, supporting its conservative reading as transdiagnostic
\emph{functional burden} rather than a symptom-liability ``\textit{p}'' factor.
Depression and anxiety instruments (MADRS, QIDS, STAI) did not form a separate factor;
they loaded $0.66$--$0.80$ on \Gfac{} as cross-loading ``windows'' onto burden. The
single prior ``biology'' candidate \emph{split} into distinct metabolic and
inflammatory factors (inter-factor correlation $\approx 0.19$), and two constructs not
in the original ten---mania and substance use---were added once their indicators were
ingested and were confirmed in the joint fit. Of the candidates, nine were confirmed
(one as a split), anhedonia was rejected, and impulsivity, negative symptoms and
sensory abnormality were declared \emph{not testable} for lack of common
indicators---each verdict an outcome of the analysis rather than an assumption
(Table~\ref{tab:dimensions}). The structure was not an artefact of the priors: a
flat-prior refit reproduced the loadings and factor correlations essentially exactly
(Tucker congruence $\varphi = 1.00$), and model comparison decisively preferred the
bifactor over unidimensional and correlated-factor alternatives. Absolute fit was
acceptable across both likelihood blocks (continuous Bayesian standardized
root-mean-square residual $\approx 0.07$; $21/22$ non-Gaussian indicators reproduced
within the $90\%$ posterior-predictive interval).

\begin{figure}[t]
  \centering
  \includegraphics[width=0.98\linewidth]{fig2_map.png}
  \caption{\textbf{The nine-dimension map.} (\textbf{a}) Posterior loading atlas
  (indicators $\times$ factors); colour encodes the absolute posterior loading, with
  rows grouped by home factor. Each specific axis is anchored by a clinically coherent
  indicator block, while \Gfac{} (functional burden) carries the cross-loading
  depression/anxiety windows. (\textbf{b}) The inter-factor correlation matrix is weak
  throughout (mean $|\text{off-diagonal}|\approx 0.10$); the general factor is
  orthogonal to the specific axes by construction (top row and column), and metabolic
  and inflammatory load correlate only $\approx 0.19$. Continuous backbone shown;
  mania and substance enter the joint mixed fit (Table~\ref{tab:dimensions}).}
  \label{fig:atlas}
\end{figure}

\subsection{Biology is the least severity-entangled domain (the central finding)}

Because the bifactor identification fixes \Gfac{}~\orth{}~specifics by construction,
we re-estimated the model allowing \Gfac{} to correlate freely with the specific axes.
Under this correlated-\Gfac{} model, metabolic and inflammatory load were the
\emph{least} entangled with functional burden: correlations with \Gfac{} of $0.12$
(metabolic) and $0.07$ (inflammatory), against $0.39$ for cognition and $0.42$ for
sleep (Figure~\ref{fig:biologyg}). Two patients with the same overall impairment can
therefore carry very different metabolic and immune loads: a patient's biological
burden rides largely on axes that the global clinical-severity picture does not
capture. This dissociation was robust to the obvious confounders. Adjusting each
indicator for age, sex, education and recruitment site, and then additionally for
antipsychotic exposure, did \emph{not} raise the biological correlations---if anything
it lowered them (metabolic $0.12\!\to\!0.07$, inflammatory $0.07\!\to\!0.06$), and the
biological axes remained the least entangled at every step
(Figure~\ref{fig:biologyg}; Annex~\ref{ann:measurement}). Consistent with a large
independent literature (Discussion), we state this as relative, not absolute,
independence: metabolic and inflammatory load are \emph{largely independent of a
general functional-burden axis}, with the $0.07$--$0.12$ values quantifying---at the
low end of what that literature would predict---a pattern previously described only
qualitatively. This is the property a biology-aware stratification can exploit, and it
is the load-bearing premise of everything that follows.

\begin{figure}[t]
  \centering
  \includegraphics[width=0.80\linewidth]{fig3_biology_g.png}
  \caption{\textbf{Metabolic and inflammatory load are the least burden-entangled
  domains.} Correlation of each continuous specific axis with the general
  functional-burden factor \Gfac{} (bars), unadjusted and after adjustment for
  medication, adiposity and site; the bifactor general loading $|\lambda_G|$
  (diamonds) is shown for comparison. Metabolic and inflammatory load fall in the
  ``biology zone'' ($\Phi<0.15$) and stay there under adjustment, whereas cognition
  and sleep partly track \Gfac{}.}
  \label{fig:biologyg}
\end{figure}

\subsection{The patient space is a continuum, not biotypes (the map organizes)}

On these nine coordinates, with each patient's measurement uncertainty propagated, the
patient space is a graded continuum rather than a set of discrete biological kinds. A
structure-discovery battery run over the posterior draws---cluster tendency, modality,
silhouette, the gap statistic, a Gaussian-mixture information criterion, density-based
clustering and a topological summary---returned a continuum verdict: the silhouette
peaked at only $0.14$ (below the $0.15$ separation floor), the first principal
component was unimodal (Hartigan dip $p\approx 1.0$), density-based clustering assigned
$100\%$ of points to noise, and the topological summary was a single connected
component. A single-Gaussian \emph{falsification null} made the verdict decisive: the
best partition of the real cloud separated patients no better than slicing a
structureless Gaussian blob of identical covariance (best silhouette $0.140$ for the
real data versus $0.141\pm0.002$ for the null, $z=-0.36$; the Gaussian-mixture-optimal
number of components collapsed to one in all $20$ posterior draws;
Figure~\ref{fig:continuum}a). There are, in this clinical--biological space, no
well-separated, reproducible discrete clusters.

We therefore represent the space by its geometry rather than by a partition. Diagnosis
is intermixed throughout the cloud (Figure~\ref{fig:continuum}b); overall burden forms
a smooth gradient along the principal axis (Figure~\ref{fig:continuum}c); and
inflammatory load varies along a \emph{different} direction
(Figure~\ref{fig:continuum}d)---the biology\,\orth\,\Gfac{} dissociation, made visible.
Four archetypal extremes, selected by cross-seed stability rather than by a
parsimony rule (the only stable resolutions were two and four corners; cross-seed
Tucker congruence $0.98$), give the continuum a clinically legible reading
(Figure~\ref{fig:continuum}e). They separate biology from symptoms from severity: a
\emph{biological} corner (high metabolic, inflammatory and substance load), a
\emph{low-burden} pole, a \emph{severe, non-biological} corner (high overall severity
but low biological load), and a \emph{symptom} corner (high sleep, developmental,
suicidality and mania). Each patient is a soft convex blend of all four (mean
assignment entropy $0.78$); only $39\%$ sit near a single pole. The reading is
transdiagnostic (adjusted Rand index $\approx 0$ against both cohort and the seven
DSM-5 subtypes), driven by the specific and biological axes rather than by severity
(variance explained by the coarsest tessellation: specific axes $\eta^2=0.12$ versus
$\Gfac{}$ $0.008$), not a missingness artefact (a classifier predicting region
membership from each patient's coverage pattern performed below baseline, lift
$-0.05$), and a tighter description of the coordinate cloud than diagnosis
(mixture information criterion $197{,}900$ versus $201{,}300$ for a DSM-5-constrained
model). Critically, because biological load varies continuously without a density
\emph{gap}, it never forms a clean partition boundary at any granularity; it surfaces
instead as an archetypal \emph{corner}. The biological phenotype is thus a direction
in a continuum, not a category---which is exactly why the continuous coordinates and
the archetype simplex, not any hard tessellation, are the load-bearing objects
downstream.

\begin{figure}[t]
  \centering
  \includegraphics[width=0.99\linewidth]{fig4_continuum.png}
  \caption{\textbf{The patient space is a continuum, not biotypes.} (\textbf{a}) A
  single-Gaussian falsification null: the best silhouette achievable on the real
  nine-dimensional cloud is no better than on a structureless Gaussian of matched
  covariance ($z=-0.36$); the mixture-optimal number of clusters is one in all
  posterior draws. (\textbf{b--d}) Principal-component embedding of the coordinates
  ($N=9{,}013$), coloured by diagnosis (fully intermixed), by general burden \Gfac{}
  (a smooth gradient), and by inflammatory load (a gradient in a \emph{different}
  direction). (\textbf{e}) Standardized profiles of the four archetypal extremes
  across the nine axes; the biological corner (A0) and the severe, non-biological
  corner (A2) carry comparable general burden but opposite biological profiles.}
  \label{fig:continuum}
\end{figure}

\subsection{The geometry persists over two years (the map persists)}

Scored forward onto the $12$- and $24$-month visits on the \emph{fixed} model
(observed cells, uncertainty propagated, never re-estimated; baseline coordinates
reproduced at Pearson $r\approx0.99$ through this path), the measurement remained
invariant for all five testable backbone axes (longitudinal Tucker congruence
$\varphi = 0.97$--$1.00$; inflammatory load, only partially invariant on a plain
Gaussian map, was invariant here at $\varphi=0.97$; \edfig{fig:invariance}). Coordinate
change can therefore be read as patient change rather than instrument drift. A
measurement-error decomposition then separated trait from state
(Figure~\ref{fig:traitstate}). The durable, trait-like axes were the biological and
cognitive ones (test--retest intraclass correlation: metabolic $0.91$, substance
$0.88$, cognition $0.70$, inflammatory $0.62$); the moving, state-like axes were the
symptom ones (sleep $0.47$, suicidality $0.42$, developmental risk $0.39$). General
burden was trait \emph{by rank} (intraclass correlation $0.62$): patients largely kept
their relative position even as the cohort as a whole improved on severity and symptoms
(population shift $-0.46$ on \Gfac{} and $-0.84$ on suicidality, while biology stayed
static). Archetype identity persisted---the archetype weight vector moved little over
two years (median cosine similarity $0.90$). The clinical reading is compact:
\emph{stratify on the durable biology; monitor the moving symptoms}.

\begin{figure}[t]
  \centering
  \includegraphics[width=0.98\linewidth]{fig5_persistence.png}
  \caption{\textbf{Durable biology, moving symptoms.} (\textbf{a}) Trait fraction
  (test--retest intraclass correlation, with $94\%$ credible intervals) over
  follow-up: the biological and cognitive axes are trait-like, the symptom axes
  state-like; mania is data-limited (hatched). (\textbf{b}) Individual rank stability
  versus population-level change: patients keep their biological rank (high intraclass
  correlation, near-zero population shift) while the cohort slides down the severity
  and symptom axes.}
  \label{fig:traitstate}
\end{figure}

\subsection{It predicts future functioning, modestly and at the group level (the map predicts)}

We then asked whether a baseline coordinate or phenotype predicts a two-year outcome
\emph{incrementally beyond} diagnosis, current severity and the baseline value of the
outcome, using an errors-in-variables Bayesian model that propagated each coordinate's
measurement uncertainty. The map added predictive signal for \emph{functioning} but
not for severity (which was saturated by its own baseline autoregression). The signal
was carried by the \emph{phenotype}, not by any single laboratory axis: the four
archetypes added the most held-out predictive value (expected log predictive density
gain $+59$ over the diagnosis-plus-severity-plus-baseline reference), the eight
specific axes added $+37$, and---an honest, map-specific shift from our earlier
analyses---the isolated ``durable biology'' axes \emph{alone} no longer cleared the
bar ($+2$, ambiguous; Figure~\ref{fig:prognosis}b). Biology still predicts functioning,
but as a \emph{corner of the phenotype} (the biological archetype) rather than as a
standalone three-axis block. No hard partition improved on the continuous and
archetype encodings (every tessellation added $\approx+20$, strictly dominated by the
archetypes); the operative number of clusters is therefore \emph{none}---the actionable
object is a continuous position, not a category. The map was \emph{co-informative} with
DSM-5 rather than redundant with it (functioning: map $+22$, DSM-5 $+29$, both $+67$;
Figure~\ref{fig:prognosis}c), and \emph{course-dependent}, with the signal concentrated
in the episodic, open-course cohorts (bipolar disorder and depression) and weak when
bipolar disorder was removed.

Translated to the clinic, the four archetypes sorted patients into groups whose
two-year functional-remission rate ranged from $27\%$ to $60\%$, transdiagnostically
(Figure~\ref{fig:prognosis}a); the \emph{biological} corner carried the worst
functional prognosis ($27\%$), worse than the equally severe non-biological corner
($32\%$)---the dissociation translating into prognosis. This gradient was a genuine
within-diagnosis effect, not a cohort-composition artefact: the rank held inside every
cohort (a biological-vs-low-burden odds ratio $\approx 4$ in each), standardizing to a
common cohort mix moved it by only $\approx 6\%$, and the corner-by-cohort interaction
was non-significant. The individual-level discrimination gain was nonetheless small
(functional-remission AUC $+0.011$ over a strong diagnosis-plus-severity reference):
the map's value is group-level prognostic stratification and continuous functional
forecasting, not a large boost to individual binary prediction. The functioning signal
survived inverse-probability weighting for attrition ($+53$), restriction to either
two of the three cohorts ($+56$, $+59$), and a permutation null (which correctly
vanished). Finally, against the full set of $143$ raw indicators under a matched
gradient-boosting classifier, the nine-dimension map was a \emph{sufficient}
representation for predicting deterioration (equal AUC) and \emph{near-sufficient} for
recovery (raw indicators $+0.04$ AUC), with $97\%$ of the raw recovery signal already
contained within the nine factors (\edfig{fig:repbench})---the compression is faithful,
not lossy in any axis the map omits.

\begin{figure}[t]
  \centering
  \includegraphics[width=0.99\linewidth]{fig6_prognosis.png}
  \caption{\textbf{Prognostic reach for functioning.} (\textbf{a}) Archetype
  prognostic atlas: two-year functional-remission rate by baseline archetype; the
  biological corner is worst ($27\%$) and the low-burden pole best ($60\%$), a
  transdiagnostic $27\%$--$60\%$ gradient. (\textbf{b}) Held-out incremental value for
  functioning: the continuous archetype encoding dominates the eight specific axes, any
  hard tessellation, and the isolated biology axes (operative number of clusters
  $=$ none). (\textbf{c}) The map and DSM-5 are co-informative---each adds beyond the
  other, and both together add most---but the individual-level remission AUC gain is
  only $+0.011$.}
  \label{fig:prognosis}
\end{figure}

\subsection{It does not yet guide treatment---a boundary drawn by the evidence}

Treatment exposures, recovered from the per-cohort source dictionaries and harmonized
to drug-class indicators, were run through an identification-first causal pipeline
(common-support overlap gate $\rightarrow$ propensity weighting $\rightarrow$
doubly-robust errors-in-variables moderation with an E-value sensitivity analysis and a
minimum-detectable-effect calculation) to ask whether the durable axes or the
archetypes \emph{moderate} treatment response. On observational, expert-centre
treatment-as-usual they did not reliably do so (\edfig{fig:treatment}a). Where
identification was cleanest---lithium in bipolar disorder, near-complete overlap---the
result was a \emph{well-identified, bounded null}: the average effect was
indistinguishable from zero (E-value $1.06$) and the design could have resolved an
interaction as small as $\approx 0.2$ standard deviations but did not. An
antipsychotic-by-biology signal for functioning in bipolar disorder was a confounded
\emph{average} effect (E-value $1.77$) with \emph{no} reliable moderation (the two map
encodings disagreed on which axis carried it---the signature of a false positive); and
clozapine in schizophrenia was channelled to treatment-resistant patients and hence
underpowered rather than cleanly null. Average treatment effects were uniformly
confounding-fragile (E-values $1.1$--$1.8$), the expected signature of confounding by
indication.

This boundary is earned, not assumed, and it cuts two ways that strengthen the rest of
the paper. First, the prognostic finding was \emph{not} merely unmodelled treatment:
the archetype carrier of the functioning forecast survived adjustment for the drug
classes patients were actually receiving (attenuation $4.7\%$; $3.9\%$ under attrition
weighting), so the biological phenotype is not a treatment proxy. Second, even without
\emph{selecting} a drug, the map \emph{describes} who faces a difficult course: the
baseline archetype stratified two-year treatment course, with the biological corner
carrying roughly twice the treatment-resistance and side-effect burden of the
low-burden corner (resistance $43\%$ vs $20\%$; significant side-effects $23\%$ vs
$11\%$), beyond baseline severity, substance comorbidity and demographics, and within
cohort (\edfig{fig:treatment}b). As in the prognostic analysis, this is a group-level
monitoring signal, not an individual classifier (the steepest gradient,
treatment-resistance, was AUC-marginal under a permutation null). Demonstrating
treatment \emph{selection} from this map would require randomized or trial-arm data,
which this baseline observational cohort---patients receiving usual care in French
expert centres---does not contain.
